\section{System Overview}
\label{sec:system-overview}

\system\ is organized around four layers: target adapters, campaign
orchestration, analysis/reporting, and the new vulnerability engine. The code
path is spread across \path{core/}, \path{harness/}, \path{analysis/}, and
\path{engine/}. The system view in \cref{fig:architecture} highlights the main
data flow.

\begin{figure*}[t]
  \centering
  \resizebox{\textwidth}{!}{\begin{tikzpicture}[
  node distance=8mm and 9mm,
  font=\scriptsize,
  box/.style={rounded corners=3pt, draw=#1, fill=#1!8, thick, align=left, inner sep=5pt, text width=0.20\textwidth},
  arrow/.style={-{Latex[length=2.5mm]}, thick, draw=raGray}
]
\node[box=raBlue] (registry) {
  \textbf{Inputs and Adapters}\\[2pt]
  Model registry\\
  Provider backends\\
  ATT\&CK-style taxonomy\\
  Probe builders
};

\node[box=raTeal, right=of registry] (harness) {
  \textbf{Execution Layer}\\[2pt]
  Campaign runner\\
  Comparison runner\\
  Trajectory store\\
  Async retries and limits
};

\node[box=raGold, right=of harness] (engine) {
  \textbf{Vulnerability Engine}\\[2pt]
  Recon phase\\
  Focused attack phase\\
  Confirmation phase\\
  Disclosure routing
};

\node[box=raRed, below=of harness] (analysis) {
  \textbf{Evaluation and Reporting}\\[2pt]
  \sr\ evaluator\\
  CVSS-like vulnerability scoring\\
  Heatmaps and comparisons\\
  Publication bundle generation
};

\draw[arrow] (registry) -- (harness);
\draw[arrow] (harness) -- (engine);
\draw[arrow] (harness) -- (analysis);
\draw[arrow] (engine) -- (analysis);
\draw[arrow] (analysis.west) -| ($(registry.south)+(0,-0.6)$) |- (registry.south);
\end{tikzpicture}
}
  \caption{High-level architecture of \system. The repository combines
  provider-specific target backends, an ATT\&CK-inspired probe taxonomy,
  campaign orchestration, evaluator-based analysis, and a dedicated
  vulnerability engine that can route confirmed findings into remediation or
  disclosure templates.}
  \label{fig:architecture}
\end{figure*}

The \textbf{target layer} provides a uniform asynchronous interface over
OpenAI, Anthropic, Google Gemini, Mistral, Together, Ollama, and vLLM-style
endpoints. These adapters normalize conversation handling, retries, rate
limits, and JSONL trajectory logging, which lets the rest of the code operate
over a common \texttt{BaseTarget} abstraction.

The \textbf{campaign layer} supplies the curated registry used in this paper
(11 aliases across OpenAI, Anthropic, Google, and Mistral), a 20-technique
ATT\&CK-inspired taxonomy, and runners for single-model and multi-model
campaigns. The stable empirical results analyzed here come from this layer:
the archived comparison run used the existing full campaign path rather than
the new engine's recon/attack/confirm execution loop.

The \textbf{analysis layer} scores campaign outputs, aggregates CVSS-like
severity estimates, maps results back to tactic categories, and produces
publication-oriented artifacts such as heatmaps, comparison tables, and LaTeX
reports. This is also where the \sr\ evaluator integration lives.

Finally, the \textbf{engine layer} introduces a dedicated vulnerability workflow
that is more focused than the legacy campaign runner. It adds explicit recon,
attack-surface ranking, confirmation, batch execution across models, and
finding routing. Because the engine is now part of the codebase even though the
archived empirical run predates a full live engine sweep, we describe it as a
systems contribution in its own right and separate it from the specific results
section below.
